\section{人类全局：既单极又多元}

上一章讨论 interface 如何决定智能边界，本章收束到平台竞争和人类全局。姚顺雨认为 OpenAI 可能成为类似 Google 的重要公司，但这不代表世界会被单极系统垄断。模型、算力、数据和 Super App 会集中，这是单极趋势；但不同 interface、不同任务、不同组织结构和 different bet 又会制造多元机会。

\begin{figure}[H]
\centering
\includegraphics[width=0.94\textwidth]{single-pole-vs-plural.png}
\caption{既单极又多元：重要模型公司会成为关键一环，但世界不必被单极垄断。自制概念图，依据 01:35:00--02:31:20 对谈内容整理。}
\end{figure}

\begin{knowledgebox}{读图：单极和多元可以同时成立}
基础模型和超级入口可能高度集中；但智能系统进入世界的方式很多，任务、界面、工具、行业和组织差异会不断制造新入口。因此“有强中心”和“有多元生态”并不矛盾。
\end{knowledgebox}

\begin{knowledgebox}{500 亿美金 allocate 问题}
\begin{tabularx}{\textwidth}{p{0.20\textwidth}p{0.34\textwidth}X}
\toprule
下注方向 & 可能买到什么 & 风险 \\
\midrule
Frontier model & 模型能力和算力规模 & 与现有巨头正面竞争，资本密度极高。 \\
Super App & 用户入口和数据回流 & 容易被已有强入口牵引，路径依赖强。 \\
Different interface & 新任务、新交互、新生态 & 不确定性高，但上限可能更高。 \\
Public-good research & 安全、评测、开放环境和工具 & 商业回报慢，但人类贡献更直接。 \\
\bottomrule
\end{tabularx}
\end{knowledgebox}

\subsection{Different Bet：超越霸主需要不同下注}

前面说世界既有集中趋势也有多元机会，本节解释多元机会从哪里来。访谈开头和结尾都反复出现 different bet。姚顺雨举例说，如果 OpenAI 一直做强化学习，也很难超过 DeepMind；要超越霸主，通常要下注到不同路径上。今天的创业机会也类似：如果只复制 ChatGPT 或 Claude，很难；如果设计不同 Super App、不同 interface、不同任务环境，就仍有空间。

\begin{figure}[H]
\centering
\includegraphics[width=0.96\textwidth]{different-bet.png}
\caption{Different Bet：超越霸主需要不同下注，而不是复制已有主线。自制概念图，依据 00:00:11--00:01:23 与 02:30:20--02:30:37 对谈内容整理。}
\end{figure}

\begin{importantbox}{创业机会的判断}
真正的危险不是“一个类似微信的东西打败微信”，而是一个不一样的东西打败微信。AI 时代的新机会也更可能来自新任务、新界面和新组织，而不是复制已有聊天产品。
\end{importantbox}

\begin{knowledgebox}{Different Bet 的判断清单}
一个下注是否真的 different，可以问四个问题：它是否改变了用户入口？是否改变了任务环境？是否改变了反馈和 reward？是否改变了组织结构或商业模式？如果只是把同样聊天框换个皮肤，它通常不是 different bet。
\end{knowledgebox}

\subsection{本章小结}

“既单极又多元”是这期最重要的平台判断。强模型公司会成为关键节点，但智能系统最终进入世界的边界，仍由不同 interface 和 different bet 决定。

